
\clearpage
\appendix

\setlength{\abovedisplayskip}{\baselineskip}
\setlength{\belowdisplayskip}{\baselineskip}
% \setcounter{page}{1}
\renewcommand{\thesection}{\Alph{section}}
% \renewcommand{\thefigure}{\Alph{figure}}
% \renewcommand{\thetable}{\Alph{table}}
\setcounter{section}{0}
% \setcounter{figure}{0}
% \setcounter{table}{0}
\setcounter{footnote}{0}

\onecolumn

\section{Appendix}

\section{Dataset}
\label{sec:A}

\noindent \textbf{Preprocessing.} All 3D images were resampled to a uniform resolution of $1\times1\times1 \text{mm}^3$. Then, 2D slices centered on the region of interest (ROI) were extracted from the axial, sagittal, and coronal views. Each slice was cropped to $192\times192$ pixels and if an image was smaller, zero padding was applied to reach the required size.  From the center of the ROI, 3 to 5 2D slices were extracted.

\vspace{3pt}
\noindent \textbf{BraTS2021} \cite{baid2021rsna} includes multimodal MRI data from 1251 glioma patients. It  has four imaging modalities: T1-weighted, T1 contrast-enhanced, T2-weighted, and FLAIR images, with image dimensions of $240\times240\times155$ and a voxel spacing of $1\times1\times1 \text{mm}^3$. The ROIs included edema, enhancing tumor, and tumor core. Of the 1251 subjects, 1141 were used for training and 110 for evaluation.

\vspace{3pt}
\noindent \textbf{BraTS2023-MEN} \cite{labella2023asnrmiccai} contains multimodal MRI scans from 1000 meningioma patients. It includes T1-weighted, T1 contrast-enhanced, T2-weighted, and FLAIR images, with dimensions of $240\times240\times155$ and a voxel spacing of $1\times1\times1 \text{mm}^3$. The ROIs include hyperintensity, enhancing tumor, and tumor core. Out of the 1000 subjects, 909 were used for training and 91 for evaluation.

\vspace{3pt}
\noindent \textbf{IXI} \footnote{https://brain-development.org/ixi-dataset/} consists of normal brain MR images. A total of 578 subjects were used in this study. It includes three modalities: T1-weighted, T2-weighted, and Proton Density, with image dimensions of $256\times256\times150$ and a voxel spacing of $0.937\times0.9375\times1.2 \text{mm}^3$. Out of the 578 subjects, 527 were used for training and 51 for evaluation.

\vspace{3pt}
\noindent \textbf{ATLAS 2.0} \cite{liew2022large} contains brain MRI data from stroke patients, including T1-weighted images and manually segmented lesion masks from 496 subjects. The T1-weighted MRI data were acquired using 1.5 Tesla and 3 Tesla MR scanners and all images have a voxel spacing of $1\times1\times1 \text{mm}^3$.

\vspace{3pt}
\noindent \textbf{UPENN-GBM} \cite{bakas2021multi} includes multimodal MR images and corresponding ROIs. A total of 147 glioma subjects were used in this study. It includes T1-weighted, T1 contrast-enhanced, T2-weighted, and FLAIR images with dimensions of $240\times240\times155$ and a voxel spacing of $1\times1\times1 \text{mm}^3$. The ROIs encompass edema, enhancing tumor, and tumor core regions.

\section{Implementation Details.}
\label{sec:B}

\noindent \textbf{Base Language Model.} Our base language model is \texttt{dolly-v2-3b} \cite{dolly2023introducing}, which originally contained 50,821 token types ($K_\texttt{text}$). We have expanded its token embedding table to 51,845 entries by adding 1,024 new image tokens ($K_\texttt{img}$). For each image token generated by the VQ-GAN \cite{esser2021taming} encoder, we add $K_\texttt{text}$ to its value and pass it to the language model as a token ID. If the model outputs a token with an ID greater than or equal to $K_\texttt{text}$, it is treated as an image token, in which case we subtract $K_\texttt{text}$ from the token ID before sending it to the VQ-GAN decoder. The model was trained using the AdamW \cite{loshchilov2018decoupled} optimizer with a batch size of 2 and a learning rate of $5\times10^{-6}$

\vspace{3pt}
\noindent \textbf{VQ-GAN.} We used VQ-GAN as our image tokenizer and trained it on $192\times192$ images. The codebook contains 1024 indices and each embedding has 256 dimensions. The encoder and quantizer converted a $192\times192$ image into a $16\times16$ matrix, which then flattened the results in 144 image tokens. The model was trained using the Adam \cite{KingmaB14} optimizer with a batch size of 2 and a learning rate of $4.5\times10^{-6}$

\vspace{3pt}
\noindent \textbf{Skip connection.} We added a skip connection block to the existing VQ-GAN and kept the VQ-GAN frozen during training. The skip connection block used inputs from all feature maps at every scale and was made up of group normalization, SiLU \cite{elfwing2018sigmoid}, a convolution layer, and cross-attention. These skip connections were added to the decoder’s feature maps, and for segmentation, the decoder’s output was passed through an extra convolution layer before producing the final result. The model was trained using the Adam optimizer with a batch size of 2 and a learning rate of $4.5\times10^{-6}$

\section{Data Generation}
\label{sec:C}
This section describes the prompt and LLM used for data generation. The prompt was based on Biomed-VITAL \cite{cui2024biomedical}.

\vspace{3pt}
\noindent \textbf{Simple Caption.} We use the LLMs to generate text data by summarizing each image with a simple caption. An example of an image caption used is shown in Fig. \ref{figA}. During data generation, the LLM is provided with an instruction, an image, and its caption. Afterward, it generates a response accordingly.

\vspace{3pt}
\noindent \textbf{Prompt for Text Data Generation.} Similar to the prompt used in Biomed-VITAL \cite{cui2024biomedical}, a detailed prompt for generating report data is shown in Fig. \ref{figB}. For generating VQA, data are shown in Fig. \ref{figC}. This prompt is fed into Gemini 1.5 Flash-8B \cite{team2023gemini}, LLaMA 3.2 11B \cite{grattafiori2024llama3herdmodels}, and GPT 4o-mini \cite{achiam2023gpt}, and the final refined report is generated by GPT 4o-mini.

\vspace{3pt}
\noindent \textbf{Evaluation Prompt.} We rate the text data produced by the LLM on a scale from 0 to 10 using a prompt that includes multiple evaluation criteria along with GPT 4o. The evaluation prompt is shown in Fig. \ref{figD}.

\section{Instruction Template}
\label{sec:D}

\noindent The instructions used for \textbf{report generation} are illustrated in Fig. \ref{figE}.

\vspace{3pt}
\noindent The instructions used for \textbf{segmentation} are illustrated in Fig. \ref{figF}.

\vspace{3pt}
\noindent The instructions used for \textbf{image translation} are illustrated in Fig. \ref{figG}.

\clearpage

\fvset{formatcom=\rmfamily\mdseries}
\begin{figure*}[t]
\centering
\begin{minipage}{1\linewidth}
    \begin{adjustbox}{scale=0.65}
    \begin{minipage}{1.52\columnwidth}
    \begin{Verbatim}[frame=single, numbersep=5pt,
                     commandchars=\\\{\}]
    - Captured using a \{\texttt{modality}\} scan, this image shows a \{\texttt{plane}\} section of the brain and demonstrates a \{\texttt{abnormality}\}.
    - This image was obtained using a \{\texttt{modality}\} scan and presents a \{\texttt{plane}\} view of the brain. It illustrates a \{\texttt{abnormality}\}.
    - The \{\texttt{plane}\} view in this image is captured via a \{\texttt{modality}\} scan of the brain. It reveals a \{\texttt{abnormality}\}.
    - Utilizing a \{\texttt{modality}\} scan, this brain image displays a \{\texttt{plane}\} perspective. The image indicates a case of \{\texttt{abnomrality}\}.
    - This brain scan employs a \{\texttt{modality}\} modality to produce a \{\texttt{plane}\} view, highlighting a \{\texttt{abnormality}\}.
    \end{Verbatim}
    \end{minipage}
    \end{adjustbox}
    \centering
    \vspace{-6pt}
    \caption{The formats of simple caption for LLM to generate text data.}
    \label{figA}
\end{minipage}

\vspace{12pt}

\centering
\begin{minipage}{1\linewidth}
    \centering
    \begin{adjustbox}{scale=0.825}
    \begin{minipage}{1.2\columnwidth}
    \begin{Verbatim}[frame=single, numbersep=5pt,
                     commandchars=\\\{\}, breaklines=true,  
                     breaksymbolleft={}, breaksymbolright={}]
\textsf{\textcolor{blue}{messages}} = [\{\textsf{"type": "text", "text":} 
    "You are an AI radiologist specialized in biomedical topics. You are provided with a brain MR image. In some cases, you may have additional text (image context) that mentions the image. Please meticulously extract all possible visual details from the image, and when generating questions and answers, ensure to integrate and consider the provided supplementary textual information. It is crucial to highlight the connections and correlations between the textual content and the visual elements within the picture to capture the full context.

    Your task is to generate a report about the image. It is essential to thoroughly consider and reference the accompanying textual information (image caption and image context) to ensure that the answers highlight the significance of the visual details present.

    Below are the requirements for generating the questions and answers in the conversation:
    - Avoid quoting or referring to specific facts, terms, abbreviations, dates, numbers, or names, as these may reveal the conversation is based on the text information, rather than the image itself. Focus on the visual aspects of the image that can be inferred without the text information.
    - Ensure that questions are diverse and cover a range of visual aspects of the image.
    - Answer responsibly, avoiding overconfidence, and do not provide medical advice.
    - The report should be generated without any sub-items or breakdowns and should not contain any line breaks. It should also avoid using special characters, except for parentheses, commas, question marks, periods, exclamation points, etc. and the language should be English only.  
    - Extract as much key detailed information from the image as possible, and I will also provide you with some text to supplement the image.
    - Avoid hallucination and ensure that the generated report is relevant to the image and the text provided."\}]

\textsf{for \textcolor{blue}{sample} in \textcolor{blue}{fewshot\_samples}}:
    \textsf{\textcolor{blue}{messages}.append({"type": "image", "image": \textcolor{blue}{sample['image']}})}
    \textsf{\textcolor{blue}{messages}.append({"type": "text", "text": \textcolor{blue}{sample['report']}})}

\textsf{\textcolor{blue}{messages}.append({{"type": "image", "image": \textcolor{blue}{image}}})}
\textsf{\textcolor{blue}{messages}.append({{"type": "text", "text": \textcolor{blue}{caption}}})}
    \end{Verbatim}
    \end{minipage}
    \end{adjustbox}
    \centering
    \vspace{-6pt}
    \caption{The prompt for report data generation.}
    \label{figB}
\end{minipage}
\end{figure*}


\begin{figure*}[t]
\centering
\begin{minipage}{1\linewidth}
    \begin{adjustbox}{scale=0.825}
    \begin{minipage}{1.2\columnwidth}
    \begin{Verbatim}[frame=single, numbersep=5pt,
                     commandchars=\\\{\}, breaklines=true,  
                     breaksymbolleft={}, breaksymbolright={}]
\textsf{\textcolor{blue}{messages}} = [\{\textsf{"type": "text", "text":} 
    "You are an AI radiologist specialized in biomedical topics.
    You are provided with a brain MR image. In some cases, you may have additional text (image Context) that mentions the image. Please meticulously extract all possible visual details from the image, and when generating questions and answers, ensure to integrate and consider the provided supplementary textual information. It is crucial to highlight the connections and correlations between the textual content and the visual elements within the picture to capture the full context.

    Your task is to generate a conversation between a person (User) inquiring about the image and you (Assistant) responding to their questions. During this interaction, the conversation should evolve as if both the User and Assistant are observing the image together. It is essential to thoroughly consider and reference the accompanying textual information (image caption and image context) to ensure a rich and informative exchange that highlights the significance of the visual details present.

    Below are the requirements for generating the questions and answers in the conversation:
    - Avoid quoting or referring to specific facts, terms, abbreviations, dates, numbers, or names, as these may reveal the conversation is based on the text information, rather than the image itself. Focus on the visual aspects of the image that can be inferred without the text information.
    - Ensure that questions are diverse and cover a range of visual aspects of the image.
    - The conversation should encompass a minimum of 3-4 exchanges of questions and answers. You may adjust the number of rounds based on the provided image and text. For content with substantial information, employing additional questions and answers may be more appropriate to ensure thorough discussion and understanding.
    - Answer responsibly, avoiding overconfidence, and do not provide medical advice. 
    - You need to start by generating the User's questions first, not assistant's answer.
    - Extract as much key detailed information from the image as possible, and I will also provide you with some text to supplement the image.
    - Avoid hallucination and ensure that the generated report is relevant to the image and the text provided.
    
    You can use the following example Questions as a guide to generate the conversation:
    - What type of scan is used for this image?
    - What abnormality do you see in this image?
    - What is the orientation of this image?"\}]

\textsf{for \textcolor{blue}{sample} in \textcolor{blue}{fewshot\_samples}}:
    \textsf{\textcolor{blue}{messages}.append({"type": "image", "image": \textcolor{blue}{sample['image']}})}
    \textsf{\textcolor{blue}{messages}.append({"type": "text", "text": \textcolor{blue}{sample['conversations']}})}

\textsf{\textcolor{blue}{messages}.append({{"type": "image", "image": \textcolor{blue}{image}}})}
\textsf{\textcolor{blue}{messages}.append({{"type": "text", "text": \textcolor{blue}{caption}}})}
    \end{Verbatim}
    \end{minipage}
    \end{adjustbox}
    \centering
    \vspace{-6pt}
    \caption{The prompt for VQA data generation.}
    \label{figC}
\end{minipage}

\end{figure*}


\begin{figure*}[t]
\vspace{-4pt}
\centering
\begin{adjustbox}{scale=0.825}
\begin{minipage}{1.2\columnwidth}
\begin{Verbatim}[frame=single, numbersep=5pt,
                 commandchars=\\\{\}, breaklines=true,  
                 breaksymbolleft={}, breaksymbolright={}]
\textcolor{blue}{\textsf{messages}} = [\{\textsf{"type": "text", "text": }
    Assume that you are a medical expert with extensive experience in your field. Your task is to assess and score a set of question-and-answer pairs from an instruction following dataset designed for fine-tuning a medical large language model (LLM). You should give a score for each Q-A pair as you are provided with multi AI conversations. You will be provided with images, their corresponding captions, in-text mentions, and the Q&A pairs. Your scoring, ranging from 0 to 10, will evaluate the following criteria:
    - Scope and Relevance: How well does the question cover key aspects of the medical image and caption provided?
    - Value for Fine-Tuning: Is the question formulated in a way that it will add value to the fine-tuning process of the medical LLM?
    - Answer Alignment: Does the provided answer directly address the question posed?
    - Accuracy: Is the information in the answer medically accurate and correct?
    - Utility: How useful is the answer in a medical context? Does it provide actionable or insightful information?
    - Image Content Recognition and Utilization: Does the response accurately identify the content depicted in the image and effectively incorporate this information into the answer to enhance comprehension or applicability in a medical context?
    
    Please consider additional factors such as:
    - Clarity: Are both the question and the answer clearly articulated and free of ambiguity?
    - Detail and Depth: Do the answer's details contribute to a deeper understanding of the topic?
    - Medical Precision: How precisely do the question and answer reflect medical terminology and knowledge?
    
    As you review each Q&A pair, please first output a single line containing scores of each Q&A pairs, splited by blank. The first score is for the first question and its corresponding answer, and so on. After that you can give some explanations, like what are the shortcomings of the current instructions.\},
    \{\textsf{"type": "image", "image": \textcolor{blue}{image}}\},
    \{\textsf{"type": "text", "text": \textcolor{blue}{caption}\}}]
    
\textsf{for  \textcolor{blue}{vqa\_data} in  \textcolor{blue}{vqa\_datas}:}
    \textsf{ \textcolor{blue}{messages}.append}(\{\textsf{"type": "text", "text":} \textcolor{blue}{\textsf{vqa\_data}}\})
\end{Verbatim}
\end{minipage}
\end{adjustbox}
\centering
\vspace{-7pt}
\caption{The prompt for data evaluation.}
\label{figD}
\vspace{-3pt}
\end{figure*}

\begin{figure*}[t]
\centering
\begin{adjustbox}{scale=0.825}
\begin{minipage}{1.2\columnwidth}
\begin{Verbatim}[frame=single, numbersep=5pt,
                 commandchars=\\\{\}]
    - Generate free-text radiology reports for the provided brain MR images.
    - Use the provided brain MR images to create corresponding free-text radiology reports.
    - Based on the provided brain MR images, produce free-text radiology reports.
    - Create free-text radiology reports that correspond to the provided brain MR images.
    - Utilize the provided brain MR images to generate corresponding free-text radiology reports.
    - Generate free-text radiology reports in accordance with the provided brain MR images.
    - Use the provided brain MR images to create accurate free-text radiology reports.
    - Produce free-text radiology reports that match the provided brain MR images.
    - Create free-text radiology reports that are consistent with the provided brain MR images.
    - Utilize the provided brain MR images to generate comprehensive free-text radiology reports.
\end{Verbatim}
\end{minipage}
\end{adjustbox}
\centering
\vspace{-7pt}
\caption{Instruction list for report generation task.}
\label{figE}
\vspace{-3pt}
\end{figure*}

\begin{figure*}[t]
\centering
\begin{adjustbox}{scale=0.825}
\begin{minipage}{1.2\columnwidth}
\begin{Verbatim}[frame=single, numbersep=5pt,
                 commandchars=\\\{\}]
    - Generate a segmentation map for \{\texttt{roi}\} in a given brain MR image.
    - Create a segmentation of \{\texttt{roi}\} in the brain MR image that highlights key structures.
    - Produce a segmented brain MR image identifying \{\texttt{roi}\}-related regions and structures.
    - Use the provided brain MR image to create the corresponding \{\texttt{roi}\} segmentation map.
    - Generate a brain MR image segmentation that accurately delineates regions of \{\texttt{roi}\}.
    - Based on the input brain MR image, produce a detailed \{\texttt{roi}\} segmentation map.
    - Use the input brain MR image to generate a segmentation highlighting areas of \{\texttt{roi}\}.
    - Create a segmentation of \{\texttt{roi}\} in the brain MR image corresponding to the input image.
    - Produce a brain MR image segmentation map that identifies regions of \{\texttt{roi}\}.
    - Generate a \{\texttt{roi}\} segmentation map from the given brain MR image to enhance visualization.
\end{Verbatim}
\end{minipage}
\end{adjustbox}
\centering
\vspace{-7pt}
\caption{Instruction list for segmentation task.}
\label{figF}
\vspace{-3pt}
\end{figure*}

\begin{figure*}[t]
\centering
\begin{adjustbox}{scale=0.825}
\begin{minipage}{1.2\columnwidth}
\begin{Verbatim}[frame=single, numbersep=5pt,
                 commandchars=\\\{\}]
    - Transform the brain MR image from \{\texttt{source}\} to \{\texttt{target}\}, 
      ensuring the output reflects the characteristics of \{\texttt{target}\}.
    - Convert the input brain MR image from \{\texttt{source}\} into a brain MR image in \{\texttt{target}\}.
    - Use the brain MR image in \{\texttt{source}\} to generate its corresponding image in \{\texttt{target}\}.
    - Generate a brain MR image in \{\texttt{target}\} based on the input from \{\texttt{source}\}.
    - Create a \{\texttt{target}\} brain MR image that mirrors the characteristics of the \{\texttt{source}\}.
    - Produce a brain MR image in \{\texttt{target}\} using the provided \{\texttt{source}\} image.
    - Transform the input brain MR image from \{\texttt{source}\} into the corresponding image in \{\texttt{target}\}.
    - Based on the provided brain MR image in \{\texttt{source}\}, create the equivalent image in \{\texttt{target}\}.
    - Generate a brain MR image in \{\texttt{target}\} that aligns with the features of the \{\texttt{source}\} image.
    - Synthesize a brain MR image in \{\texttt{target}\} using the input image from \{\texttt{source}\}.
\end{Verbatim}
\end{minipage}
\end{adjustbox}
\centering
\vspace{-7pt}
\caption{Instruction list for image translation task.}
\label{figG}
\vspace{-4pt}
\end{figure*}


